\section{Đường thứ ba trong bản đồ attribution}
\label{sec:ch16-duong-thu-ba}

\index{diễn giải cơ chế!chỗ nó đứng trong bản đồ|(}%
Đường thứ ba nhắm thẳng vào vế phải, như mục trước vừa nói, và đó là lý do người
ta trông đợi ở nó nhiều nhất: nếu đọc được bên trong mô hình thì không phải dựng
gì và cũng không phải hỏi ai. Mục này đọc bài duy nhất trong bốn bài của chương
có bản đồ đủ rộng để nói đường ấy đứng ở đâu, và cái nó nói không phải điều
người ta trông đợi.

\index{diễn giải cơ chế!bài 18 xếp nó chứ không làm nó}%
Trước hết phải nói bài ấy là bài gì, vì nhan đề chương dễ gây hiểu nhầm. Bài của
chương~\ref{ch:khung-hop-nhat} là một bài quan điểm về attribution, không phải
một bài \tn{diễn giải cơ chế}{mechanistic interpretability}. Nó không chạy thí
nghiệm nào, không báo cáo phép đo nào, có đúng một hình và ba bảng trong 25
trang. Cái nó làm là xếp: nó định nghĩa
\tn{attribution thành phần}{component attribution} và ghi rằng nhánh ấy
\emph{là một nhánh của} diễn giải cơ chế, một lĩnh vực tìm cách hiểu mô hình
bằng cách dịch ngược các cơ chế bên trong thành thuật toán đọc
được~\autocite{p18unified}. Chữ \enquote{mechanistic} xuất hiện sáu lần trong
chín trang thân bài, ba lần trong số đó ở trang đầu, còn lại nằm trong thư mục
tham khảo. Vậy bài này là bản đồ có ô dành cho đường thứ ba, không phải người đi
đường ấy.

\index{diễn giải cơ chế!ba nhánh, một bài toán}%
Quan hệ mà bản đồ ấy khẳng định là quan hệ quan trọng nhất trong mục này. Ba bài
toán attribution mà mục~\ref{sec:ch07-ba-doi-tuong} đã dựng lại, chấm điểm cho
đặc trưng, cho điểm dữ liệu huấn luyện, và cho thành phần bên trong, không phải
ba bài toán khác nhau: chúng cùng tìm một
\tn{hàm attribution}{attribution function} gán điểm cho những đối
tượng đã chọn, và \emph{chỉ khác nhau ở chỗ chọn đối tượng
nào}~\autocite{p18unified}. Phần kết của bài nói lại điều đó bằng chữ khác, rằng
ba nhánh khác nhau chủ yếu ở góc nhìn và dùng chung ba kỹ thuật lõi là
\tn{phép nhiễu}{perturbation}, gradient và
\tn{xấp xỉ tuyến tính}{linear approximation}.

\index{xấp xỉ hàm cục bộ!khung chỉ phủ attribution đặc trưng}%
Chương~\ref{ch:khung-hop-nhat} đã ghi rằng khung hình thức duy nhất của bài,
\tn{xấp xỉ hàm cục bộ}{local function approximation}, chỉ phủ attribution đặc
trưng, còn việc mở rộng sang hai nhánh kia là một giả thuyết. Mục này cần chỗ ấy
lần nữa và cần nó chặt hơn, nên đây là ba chỗ trong bài chứng thực nó. Thứ nhất,
cấu trúc mục: mục 3 có đúng một tiểu mục hợp nhất, tiểu mục về xấp xỉ hàm cục bộ
cho attribution đặc trưng, còn hai nhánh kia không có tiểu mục tương ứng nào, và
phụ lục lặp lại đúng sự lệch ấy. Thứ hai, phạm vi bài tự nêu: tám phương pháp
attribution đặc trưng, trong đó có LIME, KernelSHAP và Integrated Gradients, là
các trường hợp riêng của khung, và bảng liệt kê chúng được ghi rõ là chép lại từ
một bài trước~\autocite{p18unified}. Thứ ba, câu mở rộng: khung \emph{có thể}
cũng áp dụng cho hai nhánh kia, \emph{ta có thể giả thuyết rằng} attribution
thành phần là phép xấp xỉ hàm trên không gian các thành phần, và \emph{nếu vậy}
thì phép xấp xỉ hàm \emph{có thể} hợp nhất cả nhánh ấy~\autocite{p18unified}.
Bốn lời dè dặt trong bốn câu, và cả đoạn nằm dưới một tiêu đề về hướng đi tiếp
chứ không dưới một tiêu đề về kết quả.

\index{kiểm định chỉ số!bài 18 nói chỗ thiếu là chỗ thiếu chung|(}%
Bài phát biểu chỗ thiếu dụng cụ đo
hai lần, một lần trong thân bài và một lần trong phụ lục, và cả hai lần đều phát
biểu nó cho \emph{cả ba nhánh cùng lúc}. Ở trang 7: những khó khăn về đánh giá
phát sinh ở cả ba loại attribution; phép đánh giá
\tn{phản thực}{counterfactual} có thể cho phép kiểm
chứng chặt chẽ hơn nhưng ràng buộc tính toán thường làm nó bất khả thi; các phép
đánh giá gắn với một nhiệm vụ cụ thể thì dễ hơn nhưng thường không khái quát
được sang bối cảnh khác; đánh giá với người, dù được coi là chuẩn vàng, vướng
vấn đề mở rộng quy mô và thiên lệch; và sự đa dạng của các chỉ số đánh giá cùng
với việc chúng định nghĩa \tn{mức quan trọng}{importance} khác nhau khiến một
kết quả attribution
có thể tốt theo chỉ số này mà tồi theo chỉ số kia~\autocite{p18unified}. Ở phụ
lục, trang 24: đánh giá các phương pháp attribution gặp khó khăn đáng kể
\emph{do thiếu ground truth} và do sự phức tạp của các hệ thống AI hiện đại, và
những khó khăn ấy một phần đến từ chỗ không có định nghĩa phổ quát cho mức quan
trọng lẫn cho ground truth~\autocite{p18unified}.

\index{kiểm định chỉ số!ba cách đánh giá bài 18 liệt kê}%
Ba cách đánh giá mà bài liệt kê ở cả hai chỗ là phản thực, gắn với nhiệm vụ, và
với người. Hai trong ba là hai thứ chương này đang cân: cách thứ nhất là các họ
chỉ số của chương~\ref{ch:do-do-trung-thuc}, và cách thứ ba là đường thứ hai của
mục~\ref{sec:ch16-duong-thu-hai}. Cách thứ hai thì không, và chỗ nó lệch đáng
ghi. Bài định nghĩa phép đánh giá gắn với nhiệm vụ là phép đo mức hữu dụng thực
tế của điểm số ở một nhiệm vụ phía sau, ví dụ điểm attribution đặc trưng giúp
tìm ra những thay đổi làm lật đầu ra, hay điểm attribution dữ liệu giúp phát
hiện điểm dữ liệu gán nhãn sai~\autocite{p18unified}. Đó là hữu dụng, không phải
ground truth do thiết kế nhiệm vụ ép ra, nên nó không phải đường thứ nhất của
mục~\ref{sec:ch16-bon-duong-dung-san}. Vậy bài liệt kê hai trong ba lựa chọn của
chương này cộng thêm một lựa chọn thứ ba mà chương này không đọc, và nói cả ba
đều chưa đủ.

\index{ground truth!bài 18 không nêu ground truth nào cho thành phần}%
Có một phép rà nữa mà mục này chạy và nó cho kết quả sắc hơn cả hai đoạn trên.
Cụm \enquote{ground truth} xuất hiện đúng ba lần trong 25 trang. Hai lần đầu ở
ngay câu vừa dẫn, cả hai đều nói rằng nó thiếu: một lần cho chính chỗ thiếu, một
lần cho việc không có định nghĩa phổ quát cho nó. Lần thứ ba ở phụ lục trang 19,
và nó nói
về \tn{attribution dữ liệu}{data attribution}: các phương pháp huấn luyện lại
toàn bộ như leave-one-out thường được dùng làm một ground truth để đánh giá
những phương pháp attribution dữ liệu khác~\autocite{p18unified}. Nghĩa là trong
ba nhánh, bài nêu được ground truth cho đúng một nhánh, và nhánh ấy không phải
nhánh của đường thứ ba. Sách rà cả 25 trang và \emph{không tìm thấy một ground
truth nào được nêu cho attribution thành phần}. Nhánh mà nhan đề chương gọi là
lối thoát khỏi vấn đề ground truth, trong bài duy nhất vẽ bản đồ cả ba nhánh, là
nhánh không có ground truth nào cả.

\index{kiểm định chỉ số!bài 18 nói chỗ thiếu là chỗ thiếu chung|)}%
Đọc theo bảng~\ref{tab:ch01-bon-quan-he} thì phải ghi thêm một điều để khỏi nói
quá. Chữ \enquote{faithful} không xuất hiện lần nào trong cả 25 trang, và chữ
gần nghĩa nhất là \enquote{fidelity}, năm lần, tất cả ở phụ lục và tất cả đều
gọi tên họ chỉ số phản thực của chương~\ref{ch:do-do-trung-thuc} chứ không gọi
tên tính chất của định nghĩa~\ref{def:ch01-do-trung-thuc}. Vậy bài không nói
đường thứ ba không trung thực. Nó nói đường thứ ba không có dụng cụ đo, và nó
nói vậy về cả ba nhánh cùng một lúc.

\index{vá kích hoạt!lựa chọn tự do lặp lại ở mức thành phần}%
Chỗ thiếu ấy còn có một hình dạng quen.
Mục~\ref{sec:ch07-attribution-thanh-phan} đã đọc
\tn{phân tích trung gian nhân quả}{causal mediation analysis} và ghi rằng chữ
\enquote{làm hỏng} trong thủ tục ba lượt che một danh sách ít nhất năm lựa chọn.
Phụ lục của bài nói thêm một điều mà chương~\ref{ch:khung-hop-nhat} chưa cần
tới: trong mọi trường hợp, bộ dữ liệu dùng để sinh kích hoạt phải được chọn sao
cho nó thật sự gợi ra hành vi cần xét, kèm một chỉ số đo mức thành công của hành
vi ấy~\autocite{p18unified}. Đọc câu ấy cho kỹ thì nó đòi hai thứ có trước: phải
biết đang tìm hành vi nào, và phải có sẵn một chỉ số nói được rằng đã tìm thấy.
Đó đúng là hai lựa chọn tự do đầu trong ba lựa chọn mà
mục~\ref{sec:ch08-vong-lap-chung} tách ra cho vòng lặp phản thực, chuyển từ đặc
trưng sang thành phần. Bài không gọi chúng là lựa chọn tự do và không gọi chúng
là ground truth; cách đọc ấy là của sách.

\index{diễn giải cơ chế!chỗ nó đứng trong bản đồ|)}%
Còn một chi tiết về chi phí đáng nói vì nó bác một cách đọc dễ dãi. Nhánh tốn
kém nhất trong ba nhánh, theo lời bài, không phải nhánh thành phần: gánh nặng
tính toán nặng hơn ở các phương pháp attribution dữ liệu, vì chúng đòi huấn
luyện lại mô hình cho mỗi phép nhiễu, trong khi attribution thành phần chủ yếu
chạy ở thời điểm suy diễn~\autocite{p18unified}. Đường thứ ba vì thế không đắt
theo nghĩa phải huấn luyện lại; nó đắt theo nghĩa có quá nhiều thứ để chấm điểm,
và bài ghi rằng người ta đã rời khỏi mức từng nơ-ron vì mức ấy bất khả thi về
tính toán trên các mô hình lớn. Cái giá của đường thứ ba, nếu đọc từ bài này,
không phải tiền và cũng không phải thời gian. Nó là chỗ thiếu dụng cụ đo, thứ mà
đường này lẽ ra phải lấp và thật ra thừa kế. Mục sau đọc bài đã đo nó.
